\documentclass[10pt,letterpaper]{article}
\usepackage[T1]{fontenc}
\usepackage{amssymb}
\usepackage{amsmath}
\usepackage[margin=0.5in]{geometry}
\usepackage{siunitx}
\usepackage{xcolor}
\title{Introduction to Quantum Electronic Devices and Applications Post Lecture 16 Exercise Problem 2}
\date{23/03/2026}
\author{Jeremy Chen}
\begin{document}
	\maketitle

\noindent In insulators, $E_{F}$ is at the middle of the band gap in equilibrium, such that the overwhelming majority of occupied levels lies below $E = E_{F}$.

\noindent The band gap energy for quartz is about $E_{g} = 9.5$ eV. What is the occupation probability of a state whose energy is the lowest possible in the next highest energy band gap (i.e $E - E_{F} = E_{G}/2$)? 

\color{violet}
Given a $E - E_{F} = 4.75$ eV, the occupation probability of the state is
\begin{gather*}
f(E) = \frac{1}{1 + \exp{\left( \frac{E - E_{F}}{k_{B}T} \right)}} = \frac{1}{1 + \exp{\left( \frac{0.56}{300 * \num{8.617e-5}} \right)}} = \num{3.909e-10}
\end{gather*}
\color{black}

\end{document}